\documentclass[11pt]{article}
\usepackage[utf8]{inputenc}
\usepackage[T1]{fontenc}
\usepackage[margin=1in]{geometry}
\usepackage{microtype}
\usepackage{xcolor}
\definecolor{link}{RGB}{14,123,109}
\usepackage[colorlinks=true,urlcolor=link,linkcolor=link]{hyperref}
\pagestyle{empty}
\setlength{\parindent}{0pt}

\begin{document}
\thispagestyle{empty}

\begin{center}
{\large\bfseries Confining a Language Model to a Prescribed Role:}\\[2pt]
{\bfseries An Adversarial Reliability Test for Instruction Chatbots in Economics Experiments}\\[4pt]
{Can \c{C}elebi \quad\textperiodcentered\quad \small Preliminary, \today}
\end{center}
\vspace{0.6em}

In many economics experiments a research assistant answers participants' questions about the task. The
assistant must explain the rules, but must never give advice or hint at the best choice---that would
change the very behavior being studied. We ask a simple question: can a language-model chatbot be
trusted to stay inside this role, and what makes it reliable? We build such a chatbot for Oprea's (2024)
experiments on choice under risk, where the central debate is about separating what participants
\emph{understand} from how \emph{complex} the task is to think through. Our chatbot is therefore allowed
to clarify the rules, but is forbidden to give advice, work out the best option, or make the task look
simpler than it is. To test it without human participants, we let adversarial ``synthetic
subjects''---language models playing seven kinds of difficult participant---try to pull the chatbot out
of its role over many 15-turn conversations, more than 27,000 replies in total. An independent model
then judges every reply: did it break the no-advice rule, say something false, or fail to make sense?
Using this setup we compare different chatbot designs.

\medskip
The main result is clear. A \emph{simpler} design is safer: a single, well-instructed model call is the
most reliable, and every extra ``guardrail'' agent we add makes things worse, not better, because the
added steps tend to repeat the participant's forbidden question back to them. The one change that helps
is a \emph{stronger base model}: it cut rule-breaking replies from about 12\% to 4\% on the hardest
test, while also being more accurate and faster. Adding reasoning steps or a correction pass did not
help. We measure reliability over a whole conversation, not a single reply, because small per-reply
risks add up: even a chatbot that slips only 4\% of the time will slip at least once in most 15-turn
conversations. These results are still preliminary, but the lesson is already clear---for a chatbot that
must stay in a strict role, the quality of the base model matters far more than how elaborate the
surrounding workflow is. Full methods and interactive results:
\href{https://can-celebi.github.io/chatbot-for-experiments/}{\texttt{can-celebi.github.io/chatbot-for-experiments}}.

\end{document}
